\section{Contributions}\label{sec:intro-contributions}

In addressing the problem above, this thesis makes three contributions, revisited with supporting evidence in
Chapter~\ref{ch:conclusion}.

\begin{itemize}
  \item \textbf{A language and offline compiler.}
    We design Motivo, a typed, imperative source language for symbolic composition, and implement \texttt{motivoc},
    a staged native compiler that validates compositional structure and expands it deterministically into a
    Standard MIDI File, with no language runtime at playback.
    Version~2.0 introduces explicit declarations and signature-based overload resolution over the version~1.0
    design.

  \item \textbf{A deployable studio and engineering pipeline.}
    We build Motivo Studio, a browser-based environment that wraps \texttt{motivoc} with diagnostics, piano-roll
    preview, playback, and authenticated file storage, maintained as a tested, continuously integrated product
    rather than a one-off prototype.

  \item \textbf{An empirical evaluation.}
    We measure the compiler on synthetic benchmarks, showing that compact source expands into large MIDI timelines
    in low-millisecond compile times, with explicit typing in version~2.0 adding negligible overhead while producing
    byte-identical output.
\end{itemize}
